\documentclass[10pt,a4paper]{article}
\usepackage[margin=1.7cm]{geometry}
\usepackage[compact]{titlesec}\titlespacing*{\section}{0pt}{6pt}{2pt}
\usepackage{enumitem}\setlist{nosep,leftmargin=1.3em}
\usepackage{amsmath,amssymb,booktabs}
\usepackage[T1]{fontenc}
\usepackage{lmodern}
\usepackage{microtype}
\setlength{\parskip}{3pt}\setlength{\parindent}{0pt}
\title{\vspace{-1.2cm}Finding anapoles from the scattering matrix}
\author{Technical note, anapole study (pysie2d), 27 October 2026}
\date{\vspace{-0.8cm}}
\begin{document}
\maketitle

\section*{1. The question}
A particle is \emph{dark} for a given illumination when it does not scatter it: all the light that
comes in leaves exactly as if the particle were not there. Such a state is called an \emph{anapole}.
We want to know, for a particle of a given shape and size, at which material permittivity
$\varepsilon$ this can happen. The method below answers this using only a standard scattering solver.

\section*{2. Describing the light around the particle}
Outside the particle, any field can be written as a sum of simple wave patterns, one for each
integer order $n = 0, \pm1, \pm2,\dots$ Order $n$ varies around the particle as $e^{in\theta}$
($\theta$ is the angle around the particle centre), so $n$ counts how many times the pattern
oscillates around one turn. Each pattern comes in two versions:
\begin{itemize}
\item an \textbf{incoming} wave, travelling towards the particle, with amplitude $a_n$;
\item an \textbf{outgoing} wave, travelling away from it, with amplitude $b_n$.
\end{itemize}
We keep orders $|n|\le N$ (here $N=9$). Higher orders barely touch a particle of this size,
so they can be dropped. The amplitudes then form two vectors, $\mathbf a$ and $\mathbf b$, each
of length $2N+1$.

\section*{3. The scattering matrix $S$}
The particle is linear, so the outgoing waves depend linearly on the incoming ones:
\[
\mathbf b = S\,\mathbf a .
\]
$S$ is a square matrix, here $19\times19$. Entry $S_{mn}$ is how much outgoing order $m$ is
produced by a unit incoming wave of order $n$. If the particle were absent, every incoming wave
would simply pass through and come out unchanged, so $S$ would be the identity matrix $I$.

\textbf{How we compute it.} A solver gives the \emph{scattered} field (total minus incident)
for any illumination. We illuminate with one standing-wave pattern of order $n$ at a time and
record the scattered amplitudes of every order $m$. Those numbers form column $n$ of the
\emph{scattered-wave matrix} $T$. Because a standing wave is half incoming and half outgoing,
\[
S = I + 2T .
\]
$2N+1 = 19$ solver runs give the full matrix for one value of $\varepsilon$.

\section*{4. Why $S$ is unitary}
For a lossless particle (real $\varepsilon$), energy is conserved: the power going out equals
the power coming in, for \emph{every} incoming vector $\mathbf a$. In matrix terms
$\|\mathbf b\|^2 = \|\mathbf a\|^2$, which holds for every $\mathbf a$ only if
\[
S^{\mathsf H} S = I \qquad (S^{\mathsf H}\text{ is the conjugate transpose}),
\]
i.e.\ $S$ is \emph{unitary}. This gives a built-in check: we compute $\max|S^{\mathsf H}S - I|$
and require it to be at round-off level. Measured: $2\times10^{-13}$ for a circle,
$\le 10^{-9}$ for the non-circular shapes.

\section*{5. Channels and phase shifts}
A unitary matrix has a complete set of orthogonal eigenvectors $\mathbf v_j$, and every eigenvalue
has magnitude one, so it can be written as a pure phase:
\[
S\,\mathbf v_j = e^{2i\delta_j}\,\mathbf v_j .
\]
\begin{itemize}
\item $\mathbf v_j$ is a \textbf{channel}: a particular mix of incoming orders that comes back out
as the same mix. The particle can only delay it.
\item $\delta_j$ is the channel's \textbf{phase shift}, the delay the particle imposes on it.
\end{itemize}
The scattered part of channel $j$ is $(e^{2i\delta_j}-1)/2$, whose size is $|\sin\delta_j|$. So:
\begin{itemize}
\item $\delta_j = 0$ (or a multiple of $180^\circ$): the channel passes untouched. \textbf{This is an anapole}
for that channel.
\item $\delta_j = \pm 90^\circ$: the channel scatters as strongly as possible. This is a \textbf{resonance}.
\end{itemize}

\section*{6. Sweeping the permittivity}
As $\varepsilon$ increases, the particle stores more energy, and each phase shift moves steadily
in one direction (it never turns back). An anapole is therefore a point where one channel's
$\delta_j$ crosses zero. Between two anapoles of the same channel, its phase passes through
$\pm90^\circ$, a resonance. So anapoles and resonances alternate.

The procedure is:
\begin{enumerate}
\item Compute $S(\varepsilon)$ on a fine grid (step $0.02$ from $\varepsilon = 1.5$ to $25$).
\item Find the eigenvalues and eigenvectors at each step.
\item Follow each channel from one step to the next by picking the eigenvector most similar to the
previous one. Sorting by phase alone fails: the dropped high orders all sit at $\delta\approx0$
and get confused with a real crossing.
\item Record where a followed $\delta_j$ changes sign. Check it by solving again at the
interpolated point. A very sharp resonance can jump from $+90^\circ$ to $-90^\circ$ within one step
and look like a sign change. At a true anapole, $\delta_j$ stays small there.
\end{enumerate}

\textbf{Symmetry.} Our shapes are mirror-symmetric (left--right). Every channel is then either
even or odd under the mirror, and the two families never mix. We split $S$ into an even block and
an odd block and treat each separately, which keeps the eigenvectors clean.

\section*{7. A given illumination}
Real illumination (a plane wave, or the evanescent tail of a waveguide mode) is a fixed incoming
vector $\mathbf a$. Writing it in channels, the scattered power is
\[
P_{\mathrm{sca}} \;\propto\; \sum_j \bigl|\mathbf v_j^{\mathsf H}\mathbf a\bigr|^2 \,\sin^2\delta_j .
\]
\begin{itemize}
\item $|\mathbf v_j^{\mathsf H}\mathbf a|^2$ is how much of the illumination enters channel $j$.
\item $\sin^2\delta_j$ is how strongly that channel scatters.
\end{itemize}
The particle is dark for $\mathbf a$ only if every channel that $\mathbf a$ feeds has a small
phase shift. On a circle each order is its own channel, so one exact anapole can make one order
exactly dark. A non-circular shape mixes orders, so a fixed illumination generally feeds several
channels. It is then only \emph{nearly} dark: the sum is small but not zero.

\section*{8. Checks and results}
\begin{tabular}{@{}ll@{}}
\toprule
check & result \\
\midrule
Unitarity, $\max|S^{\mathsf H}S-I|$ & $2\times10^{-13}$ (circle), $\le10^{-9}$ (shapes) \\
Circle anapoles vs.\ exact formula & all 9 recovered: 5.535, 7.053, 10.287, 14.288, 14.320,\\
 & \quad 17.869, 18.941, 23.564, 24.206 \\
Dark design vs.\ direct simulation & $\times316$ darker than the circle, both ways \\
\bottomrule
\end{tabular}

\smallskip
\textbf{Findings} (silicon, $\varepsilon = 12.08$, equivalent radius 500\,nm, $\lambda = 1550$\,nm):
\begin{itemize}
\item Neither optimised dark shape has an anapole at $\varepsilon = 12.08$. An interior
calculation puts the nearest at 11.95 and 14.15. Off the circle, these are points where a
channel's $\delta_j$ comes close to zero (0.05, 0.015) without reaching it.
Following channels through near-crossings is fragile: on these shapes it produced false
crossings, so non-circular roots must be confirmed by an interior calculation.
\item The design that is dark above the waveguide is only $1.7\times$ darker than a circle when
the same evanescent wave hits it in free space, against $27\times$ above the slab. Light
reflected by the slab is part of the cancellation.
\end{itemize}

\textbf{Limits.} The method needs a lossless particle (otherwise $S$ is not unitary and the
phases are not pure). It describes the particle alone in free space. Results are from one solver
(pysie2d) and are not externally validated.
\end{document}
